\documentclass[11pt,a4paper]{article}
\usepackage{../../../_templates/latex/mastery-style}

\newcommand{\TermEN}[1]{\textbf{#1}}
\newcommand{\TermZH}[1]{\textcolor{MMuted}{#1}}

\begin{document}

\MasterySetBrand{25Maths}{Mastery Series}
\MasterySetHeader{Quadratics Mastery Pack}{Module 4: Vocabulary Sheet (EN + ZH)}
\MasterySetFooter{Quadratics v1 | Module 4 | Vocabulary}

\MasteryModuleCover
  {4}
  {Vocabulary Sheet (EN + ZH)}
  {Standardise method language and final-answer phrasing for exam reliability}

\section{Core Terms}
\renewcommand{\arraystretch}{1.2}
\begin{tabularx}{\linewidth}{l l X}
\toprule
\textbf{EN Term} & \textbf{ZH Label} & \textbf{Exam Meaning} \\
\midrule
\TermEN{quadratic equation} & \TermZH{er ci fang cheng} & equation with highest power 2 \\
\TermEN{root / solution} & \TermZH{gen / jie} & value making equation true \\
\TermEN{factorise} & \TermZH{yin shi fen jie} & rewrite as product of factors \\
\TermEN{discriminant} & \TermZH{pan bie shi} & \(b^2-4ac\), classifies root type \\
\TermEN{completing the square} & \TermZH{pei fang fa} & rewrite to \((x+p)^2+q\) \\
\TermEN{axis of symmetry} & \TermZH{dui cheng zhou} & vertical line through vertex \\
\TermEN{vertex} & \TermZH{ding dian} & turning point of parabola \\
\TermEN{repeated root} & \TermZH{chong gen} & same root appears twice \\
\TermEN{x-intercept} & \TermZH{x zhou jie ju} & point where graph meets x-axis \\
\bottomrule
\end{tabularx}

\section{Sentence Frames}
\begin{keyconceptbox}
\textbf{Method choice:} I selected \underline{\hspace{2cm}} because the equation is in \underline{\hspace{2.6cm}} form.
\end{keyconceptbox}

\begin{examtipbox}
\textbf{Final line format:} Therefore, the roots are \(x=\)\underline{\hspace{1.2cm}} and \(x=\)\underline{\hspace{1.2cm}}.
\end{examtipbox}

\begin{commonmistakebox}
\textbf{Context filter:} If a root is negative for a length/time question, reject it explicitly and state why.
\end{commonmistakebox}

\section{Quick Check}
Write one bilingual sentence using each frame:
\begin{enumerate}[leftmargin=1.3em]
  \item Method choice sentence (EN + ZH)
  \item Multi-root final-answer sentence (EN + ZH)
  \item Context-validity sentence (EN + ZH)
\end{enumerate}

\end{document}
